\chapter{Read this first}

This booklet explains what was built after you gave the instructions: what each piece is,
why it exists, and how the code matches the papers. It assumes you know what the project is
about and that you have heard the words ISTA and LISTA. It does not assume you remember how
any of them work, so everything is rebuilt from the start.

\begin{careful}[One correction before anything else]
You described the first network as ``the convolutional ALISTA part''. What runs is
\textbf{not} Conv ALISTA. It is \textbf{CSC-LISTA (Eq.~20 of the ALISTA paper, Liu et al.\ 2019)}.

The paper has two convolutional networks, and the names are easy to mix up:
\begin{itemize}
\item \textbf{CSC-LISTA, Eq.~20:} every filter is learned from data. This is what we use. The paper
calls it ``Conv LISTA''; we call it CSC-LISTA, for convolutional sparse coding LISTA, so that it is not
confused with the old ConvLISTA (in-house) from October that the professor criticised.
\item \textbf{Conv ALISTA, Eq.~26:} the filters are calculated by a formula, and only a few
scalars are learned. We tried a small version, you rejected it, and its code was deleted.
\end{itemize}
So ``the ALISTA paper'' is correct, but ``ALISTA'' as the name of our network is not. Say
``CSC-LISTA, Eq.~20 of Liu et al.''
\end{careful}

\section*{The journey on one page}

\begin{center}
\begin{tikzpicture}[node distance=3.5mm,
  stage/.style={draw=accent,fill=bluewash,rounded corners=2pt,text width=12.6cm,inner sep=5pt,font=\small},
  arrow/.style={-{Latex[length=2mm]},draw=ink,thick}]
\node[stage] (a) {\textbf{1. The professor's criticism.} The old ConvLISTA (in-house) had no published source, so the
project restarted at one standoff pair: 3\um\ to 0.5\um.};
\node[stage,below=of a] (b) {\textbf{2. Too few real pairs.} Only 15 COMSOL component pairs exist.
So we simulate training data and keep COMSOL as the test.};
\node[stage,below=of b] (c) {\textbf{3. Simulated data, v1 to v3.} Draw wires, compute all three field
components with physics, add COMSOL-like speckle. Chapter~\ref{ch:simulation}.};
\node[stage,below=of c] (d) {\textbf{4. Choosing a published network.} The ALISTA paper gives Eq.~20. We tried Eq.~26,
then dropped it. Chapter~\ref{ch:choice}.};
\node[stage,below=of d] (e) {\textbf{5. CSC-LISTA.} Learn a dictionary, then train 16 layers. The result was
39.1\% error on COMSOL, against 49.3\% for Wiener. Chapters~\ref{ch:csc_lista} to~\ref{ch:training}.};
\node[stage,below=of e] (f) {\textbf{6. The density problem.} The codes stopped being sparse, so we tried a threshold
floor and a relative loss. Chapter~\ref{ch:density}.};
\node[stage,below=of f] (g) {\textbf{7. Strict runs.} You asked for no additions. Strict 32 and 64 layers gave
35.4\% and 33.8\%, the best CSC-LISTA so far. Chapter~\ref{ch:results}.};
\node[stage,below=of g] (h) {\textbf{8. Two more networks.} A learned-FISTA style network scored 34.2\%, and FISTA-Net
39.4\% without converging. Chapters~\ref{ch:lfista} and~\ref{ch:fistanet}.};
\node[stage,below=of h] (i) {\textbf{9. FISTA rescored.} A wider search for $\lambda$ moved FISTA-300 to 30.5\%, and FISTA was scored
at 32 and 64 iterations for a fair comparison. Chapter~\ref{ch:results}.};
\node[stage,below=of i] (j) {\textbf{10. FISTA-Net rerun.} With the authors' start values and 9,000 updates it scored 33.3\%, the best
trained network so far. Chapter~\ref{ch:fistanet}.};
\foreach \x/\y in {a/b,b/c,c/d,d/e,e/f,f/g,g/h,h/i,i/j} \draw[arrow] (\x) -- (\y);
\end{tikzpicture}
\end{center}

\section*{The result in one sentence}

FISTA at 300 iterations, a classical algorithm with no trained layers, scores 30.5\% on COMSOL. The best trained network,
FISTA-Net, scores 33.3\%, and strict Eq.~20 at 64 layers scores 33.8\%. All of them beat the Wiener filter's 49.3\% by a wide
margin, and every score here is an error, so lower is better.

\section*{How to use this booklet}

\begin{itemize}
\item \textbf{Blue boxes} hold the idea to remember.
\item \textbf{Orange boxes} hold a trap or an honest limit.
\item \textbf{Green boxes} hold what to say if your professor pushes on that point.
\item \textbf{Grey code blocks} are copied from the repository, unchanged, apart from
\texttt{\# ...} where lines are skipped. The file name is above each block.
\end{itemize}

If you only have an hour, read Chapters~\ref{ch:problem}, \ref{ch:density}, \ref{ch:results}
and~\ref{ch:defence}.

\begin{careful}[The numbers in this edition]
All server runs have finished, so every table now holds a real score. Each run used one random seed. No number here
is a prediction.
\end{careful}
